\documentclass[UTF8,a4paper,12pt]{ctexart}

\usepackage{amsmath,amssymb,bm}
\usepackage{booktabs}
\usepackage{geometry}
\usepackage{longtable}
\usepackage{array}
\usepackage{enumitem}
\usepackage{microtype}
\usepackage{hyperref}

\geometry{left=2.6cm,right=2.6cm,top=2.5cm,bottom=2.5cm}
\setlength{\parindent}{2em}
\setlength{\parskip}{0.35em}
\renewcommand{\arraystretch}{1.35}
\setlist{nosep}
\hypersetup{hidelinks}

\title{方法与案例研究部分}
\author{}
\date{}

\begin{document}

\maketitle

\section{3 面向换乘客流汇聚的多线换乘站疏散建模}

大型多线换乘站的应急疏散不是若干条独立出站路径的简单叠加。站台候车乘客、列车到达乘客、站厅滞留乘客和跨线换乘乘客会在楼梯、扶梯、换乘通道、闸机和出口前区域等有限设施处共同竞争通行能力。与普通单线车站相比，换乘站的关键困难在于：换乘客流跨越多个线路空间，其行进链条长、涉及共享设施多，且容易与列车到达客流和站台出站客流在相近时段内汇聚。因此，全站疏散效率往往不只由最短路径长度决定，而受共享设施排队、下游接收能力和尾部清空过程共同控制。

既有地铁站疏散优化研究通常采用“疏散对象分析--仿真或网络建模--路径/控制优化--案例验证”的组织方式。Guo和Zhang（2022）在地铁站多目标疏散优化中先分析车站疏散目标，再建立仿真和代理模型，最后开展优化；Zheng和Mou（2023）与Wen等（2024）则把枢纽内部客流传播、动态加载和节点排队作为模型主体；Wei等（2026）在消防路线研究中先构建模型，再给出求解方法，并在案例中先验证节点时间预测，再验证路线效果。本文借鉴这种论文组织逻辑，但研究对象转向大型五线换乘站中的整体乘客疏散组织，核心关注换乘客流参与共享设施汇聚后对全站疏散效率的影响。

本文提出队列感知的动态路径组织模型 AdaptiveQueueAwareAStar，简称 AA。AA 不是把疏散问题简化为静态最短路，也不把对照方法解释为车站已有官方预案。其作用是在多来源客流持续进入站内网络时，预测各批次到达共享设施时将面对的队列和等待，并据此组织能够在统一容量约束下执行的疏散路径。模型由站内疏散网络、多来源客流批次、共享设施服务容量、到达时刻负荷预测和 AA 路径组织五个部分组成。

\subsection{3.1 多线换乘站疏散过程分析}

龙阳路站具有多线路换乘属性，站内疏散过程中同时存在出站、换乘和列车到达释放等不同客流来源。对于这类车站，单个乘客组的路径选择会改变共享设施的未来负荷，而共享设施未来负荷又会反过来影响后续乘客组的路径代价。若路径组织只依据当前局部拥挤程度或当前边权进行判断，就可能低估后续批次到达瓶颈设施时的排队等待，导致多个来源组在时间上集中到达同一楼扶梯、换乘通道、闸机或出口前区域。

本文将疏散过程理解为两个相互耦合的层次。第一层为路径决策层，负责为仍处于可决策阶段的客流批次确定下一段路径或完整路径。第二层为设施服务执行层，负责按照共享容量和下游空间接收约束释放实际通行人数。路径决策层回答“乘客应该朝哪里组织”，设施服务执行层回答“这些乘客能否在当前时间步实际通过设施”。只有把两层放在同一模型中，才能避免同一设施被多条路径重复计算容量，也才能把换乘客流和其他来源客流在同一设施上的竞争关系表达出来。

因此，本文的建模目标不是获得一条全站固定疏散路线，而是在离散疏散过程中持续更新以下信息：各来源组当前处于何处，哪些批次仍可调整路径，哪些批次已经承诺到达某一共享设施，关键设施未来会在何时出现排队，最终哪些路径组织能够减少全站尾部滞留和关键设施等待暴露。

\subsection{3.2 站内疏散网络与客流来源建模}

将车站公共区抽象为有向图
\[
G=(V,E),
\]
其中，节点集合 \(V\) 表示站台、站厅、换乘节点、楼梯、扶梯、闸机、出口前区域和地面出口，边集合 \(E\) 表示相邻空间或设施之间允许乘客移动的方向性连接。每条边 \(e=(i,j)\) 包含长度 \(l_e\)、有效宽度 \(w_e\)、设施类型、行进方向和容量属性；每个节点 \(v\) 包含面积 \(A_v\)、当前人数 \(N_v(t)\)、节点类型和空间接收上限 \(\bar N_v\)。

为表达多条路径对同一设施服务能力的竞争，定义边到共享设施资源的映射
\[
\rho:E\rightarrow R\cup\{\varnothing\},
\]
其中，\(R\) 为有限通行设施资源集合。若 \(\rho(e)=r\)，则边 \(e\) 的通行需要消耗共享设施 \(r\) 的服务能力；若 \(\rho(e)=\varnothing\)，则该边只表示普通空间移动。该映射用于保证同一楼梯、扶梯、闸机组或出口前区域不会因被多条边引用而重复释放容量。

客流来源组集合记为 \(\mathcal{G}\)。来源组至少包括站台等待客流、列车到达客流、站厅滞留客流和换乘客流。对换乘客流，模型保留其来源线路、目标线路或离站目标，使其在关键设施处的通过量和等待暴露能够与普通出站客流分开统计。一个进入路径选择阶段的客流批次记为
\[
b=(g,o_b,n_b,t_b),
\]
其中，\(g\in\mathcal{G}\) 为来源组，\(o_b\) 为批次当前节点，\(n_b\) 为批次人数，\(t_b\) 为批次进入路径选择阶段的时刻。采用批次建模可以保留列车到达和换乘客流释放的时间结构，也便于在路径搜索后登记其未来设施到达需求。

\subsection{3.3 共享设施服务与节点排队模型}

设共享设施 \(r\) 的服务率为 \(\mu_r\)，仿真步长为 \(\Delta t\)。时间步 \(t\) 内设施 \(r\) 可释放的整数服务能力为
\[
C_r(t)=\left\lfloor \mu_r\Delta t+\xi_r(t)\right\rfloor ,
\]
其中，\(\xi_r(t)\) 表示上一时间步未形成完整整数人的服务能力余量。若 \(A_r(t)\) 为本时间步请求进入设施 \(r\) 的人数，\(Q_r(t)\) 为设施入口等待人数，则队列更新为
\[
Q_r(t+\Delta t)=\max\{Q_r(t)+A_r(t)-C_r(t),0\}.
\]
当多个来源组在同一时间步请求同一设施时，它们共同竞争 \(C_r(t)\)。未获得服务能力的人数保留在上游等待，并计入后续时间步的队列状态。

除设施服务能力外，下游空间接收能力也会约束通行。设节点 \(v\) 当前占用人数为 \(N_v(t)\)，本时间步已被其他移动请求预留的进入人数为 \(B_v(t)\)，则新进入人数 \(y_v(t)\) 满足
\[
N_v(t)+B_v(t)+y_v(t)\leq \bar N_v .
\]
若下游节点或边的接收条件不足，即使上游设施仍有名义服务能力，乘客也不能被释放到下游空间。该约束用于刻画高负荷换乘站中常见的通道回溢和出口前区域滞留。

普通移动边的物理行走时间由边长和速度模型确定：
\[
\tau_e(t)=\frac{l_e}{v_e(t)}.
\]
其中，\(v_e(t)\) 可由速度-密度关系、设施类型或仿真设定给出。本文不把速度-密度阈值作为主要敏感性对象；其具体取值在参数表中单独列明来源，并与车站 CAD 几何、模型设定值和必要假设分开标注。

\subsection{3.4 到达时刻负荷预测模型}

AA 的关键区别在于预测“批次到达设施时”的排队状态，而不是只读取当前时刻的队列。设当前时刻为 \(t\)，共享设施 \(r\) 已登记的未来到达事件集合为
\[
\mathcal{A}_r(t)=\{(\eta_j,n_j,g_j)\},
\]
其中，\(\eta_j\) 表示已承诺批次预计到达设施 \(r\) 的时刻，\(n_j\) 表示人数，\(g_j\) 表示来源组。若候选批次 \(b\) 预计在 \(\eta\) 时刻到达设施 \(r\)，则其到达时刻队列估计为
\[
\widehat Q_r(\eta)=\mathcal{Q}\left(Q_r(t),\mathcal{A}_r(t),\mu_r,t,\eta\right).
\]
队列演化算子 \(\mathcal{Q}\) 在当前时刻到预计到达时刻之间释放服务能力，并在已登记到达事件发生时加入相应人数。若 \(\mu_r>0\)，候选批次在设施 \(r\) 处的预计等待时间为
\[
W_r(\eta)=\frac{\widehat Q_r(\eta)}{\mu_r}.
\]

未来到达事件包括两类。第一类来自已经处于执行过程的批次，其到达时间可由当前位置、边长和移动状态推算。第二类来自同一决策轮中已经完成路径搜索的前序批次。当前序批次路径被接受后，其预计进入共享设施的时间、人数和来源组立即写入 \(\mathcal{A}_r(t)\)。后续批次搜索路径时会读取这些事件，因此能够感知前序批次对未来设施容量的占用。该机制使多来源客流不再基于同一份未更新状态独立选择路径，而是在路径组织层形成可解释的设施竞争关系。

\subsection{3.5 队列感知动态路径组织模型}

对批次 \(b\) 的候选路径 \(P=(e_1,e_2,\ldots,e_m)\)，AA 将路径代价定义为
\[
J_b(P)=\sum_{i=1}^{m}\left[\tau_{e_i}(\eta_i)+W_{\rho(e_i)}(\eta_i)+S_{\rho(e_i)}(n_b)+H_{e_i}(\eta_i)\right],
\]
其中，\(\eta_i\) 为预计进入第 \(i\) 条边或设施的时刻，\(\tau_{e_i}\) 为物理行走时间，\(W_{\rho(e_i)}\) 为到达时刻设施队列等待，\(S_{\rho(e_i)}(n_b)\) 为当前批次自身消耗设施服务能力形成的服务时间，\(H_{e_i}\) 为下游空间接收不足造成的预计等待。若边 \(e_i\) 不对应共享设施，则 \(W_{\rho(e_i)}\) 与 \(S_{\rho(e_i)}\) 取 0。

路径代价按时间递推。前一段路径上的行走和等待会改变下一设施的到达时刻，而下一设施的等待又取决于这一到达时刻下已经登记的未来负荷。因此，AA 属于时间依赖路径组织方法。其改进点不是简单增加拥挤惩罚，而是把换乘客流、列车到达客流和站台客流对共享设施的未来到达需求纳入同一条路径的代价计算。

路径搜索标签记为
\[
L=(v,\eta,c,\pi),
\]
其中，\(v\) 为当前节点，\(\eta\) 为预计到达当前节点的时刻，\(c\) 为累计代价，\(\pi\) 为前驱标签。标签从节点 \(u\) 沿边 \(e=(u,v)\) 扩展时，先计算物理行走时间，再根据预计进入共享设施的时刻查询 \(\widehat Q_r(\eta)\)，最后检查下游空间接收条件，得到新的到达时刻和累计代价。由于两个标签即使到达同一节点，其到达时刻不同也会面对不同队列状态，因此不能把动态路径搜索退化为静态节点最短路。

\subsection{3.6 疏散组织效果评价方法}

评价指标分为整体效率、设施瓶颈、来源组机制和跨模型一致性四类。整体效率指标包括 \(T_{50}\)、\(T_{80}\)、\(T_{95}\)、\(T_{100}\)、平均疏散时间、累计停滞人秒、总移动距离和计算耗时。其中，\(T_{95}\) 用于观察大多数乘客清空后的尾部阶段，\(T_{100}\) 用于描述最后一批乘客离站时间，累计停滞人秒用于衡量排队、阻塞和等待暴露。

设施瓶颈指标包括关键楼扶梯、换乘通道、闸机和出口前区域的通过人数、峰值队列、队列人秒、首次排队时间、最后排队时间和设施利用率。来源组机制指标按站台等待、列车到达、站厅滞留和换乘来源统计清空时间、出口选择、关键设施通过量和等待暴露。跨模型一致性指标用于比较宏观网络模型与 Pathfinder 微观仿真的累计疏散曲线、拥堵位置和出口利用趋势。

\section{4 AA 模型求解与动态加载算法}

第 3 章给出了 AA 的模型定义，本章说明模型如何在离散疏散仿真中求解。与一次性优化全站固定路线不同，AA 在仿真过程中循环执行状态读取、批次生成、时间依赖搜索、未来事件登记、共享容量加载和指标更新。该求解过程对应许多地铁站疏散优化论文中的“模型求解”部分，但其重点是把路径组织与设施服务执行连接起来。

\subsection{4.1 计算状态与时间步推进}

设时间步 \(t\) 的系统状态为
\[
\Omega(t)=\left(G,\mathcal{B}(t),\mathcal{Q}(t),\mathcal{N}(t),\mathcal{A}(t)\right),
\]
其中，\(\mathcal{B}(t)\) 为当前可决策批次集合，\(\mathcal{Q}(t)\) 为共享设施队列状态，\(\mathcal{N}(t)\) 为节点占用状态，\(\mathcal{A}(t)\) 为未来到达事件集合。每一时间步内，模型先更新上一时间步已经完成移动的乘客，再生成本时间步自然到达的站台、列车、站厅和换乘批次。随后，AA 对可决策批次逐一搜索路径，并把被接受路径对共享设施产生的未来到达需求写回系统状态。

可决策批次按照进入路径选择阶段的时刻、来源组、当前节点和批次编号进行确定性排序。排序不是为了人为赋予某一来源组更高优先级，而是为了保证同一时间步内的事件登记具有可重复性。前序批次确定路径后，后续批次能够看到已经承诺的设施到达负荷，从而形成可解释、可复现的批次间相互影响。

\subsection{4.2 时间依赖 AA 搜索}

AA 对每个批次执行一次时间依赖 A* 搜索。搜索入口为批次当前节点，目标为可用出口集合。启发式函数采用静态自由流条件下的时间下界，用于减少搜索空间；真实路径代价仍由物理行走时间、到达时刻设施等待、当前批次服务时间和空间回溢等待共同决定。模型不缓存动态完整路径，也不预先固定 \(K\) 条候选路径，以避免在设施预测状态已经变化后复用失效的动态路径评价。

标签扩展过程包括四步。第一，根据边长和当前速度估计物理行走时间。第二，若边对应共享设施，则根据预计到达时刻查询未来队列并计算设施等待。第三，把当前批次自身服务时间加入路径代价，避免多个批次同时把同一设施服务能力视为未被占用。第四，检查下游节点或边的接收条件，估计空间阻塞等待。完成扩展后，标签的到达时刻和累计代价同步更新。

当搜索到出口节点时，AA 输出该批次的路径、预计总代价、关键设施进入时刻和代价分解。代价分解包括物理行走时间、设施等待时间、批次服务时间和空间回溢等待时间。该分解结果将在案例研究中用于说明路径改变究竟来自行走距离缩短，还是来自设施排队减少。

\subsection{4.3 未来事件登记与批次路径重评估}

批次路径被接受后，模型并不立即认为该批次已经通过相关设施，而是把其预计到达共享设施的时刻、人数和来源组登记为未来事件。事件登记的意义在于把路径决策层的承诺需求传递给后续批次，使后续批次搜索时能够感知已经形成的未来设施负荷。

对仍处于选择阶段且尚未进入边内移动或设施服务的批次，AA 允许在满足稳定性条件时重新评估路径。设保留当前路径的剩余代价为 \(C_{\mathrm{stay}}\)，候选切换路径代价为 \(C_{\mathrm{switch}}\)，相对收益为
\[
R=\frac{C_{\mathrm{stay}}-C_{\mathrm{switch}}}{C_{\mathrm{stay}}}.
\]
只有当候选路径有效、下一步方向发生变化、\(C_{\mathrm{switch}}<C_{\mathrm{stay}}\)，且 \(R\) 超过防摆动阈值时，路径切换才被接受。已经进入边内移动、已经获得设施服务能力或已经处于不可重评估状态的批次继续执行既有移动，不在途中瞬时改路。

\subsection{4.4 共享容量动态加载}

路径搜索生成的是移动意图，实际通行人数由共享容量动态加载过程决定。对同一设施 \(r\)，无论请求来自哪条线路、哪个站台、哪个换乘关系或哪条路径，只要映射到同一资源，就共同使用当前时间步的服务能力 \(C_r(t)\)。当请求人数超过服务能力时，未被释放的乘客保留在上游队列，并继续参与后续时间步竞争。

动态加载还同步检查下游空间接收条件。若下游节点达到接收上限，相关移动不会被释放到下游空间，而是记录为空间阻塞或上游等待。执行完成后，系统更新节点人数、来源组分布、设施队列、在途批次、出口统计和各类指标。该过程保证客流守恒：任一批次在同一时间步内只能处于上游等待、设施服务、边内移动、下游节点或出口中的一种状态。

\subsection{4.5 计算对照方法与实验公平性}

ImprovedAStar 在本文中仅作为计算对照方法，不代表龙阳路站运营方制定或发布的车站运营预案。该方法与 AA 使用相同车站网络、相同客流输入、相同共享容量动态加载过程、相同速度模型、相同空间接收约束和相同终止条件。两者差异限定在路径组织逻辑：AA 使用未来到达事件和到达时刻队列预测，ImprovedAStar 不显式登记未来设施负荷，也不预测候选批次到达设施时的队列等待。

\begin{table}[htbp]
\centering
\caption{AA 与计算对照方法的实验差异控制}
\small
\begin{tabular}{p{0.24\textwidth}p{0.30\textwidth}p{0.34\textwidth}}
\toprule
比较维度 & AA & ImprovedAStar \\
\midrule
路径评价依据 & 当前状态、未来到达事件、到达时刻队列和空间接收状态 & 当前网络状态和即时边权评价 \\
批次间相互影响 & 前序批次路径确定后登记未来设施负荷 & 不登记同一轮前序批次的未来设施负荷 \\
设施容量执行 & 与对照方法共用共享容量动态加载过程 & 与 AA 共用共享容量动态加载过程 \\
路径重评估 & 仅对仍处于选择阶段的批次按收益阈值重评估 & 不执行队列感知的路径重评估 \\
论文定位 & 本文提出的路径组织模型 & 计算基准，不作为车站运营预案 \\
\bottomrule
\end{tabular}
\end{table}

\section{5 龙阳路站案例研究}

本文以龙阳路站为案例验证 AA 在大型多线换乘站疏散组织中的作用。案例研究的写作顺序参照地铁站疏散优化和消防路线优化论文的常见安排：先说明案例背景和场景构造，再验证模型中间环节，随后比较整体疏散性能，最后用关键设施、换乘客流分解、Pathfinder 对照和适用边界分析解释结果。由于最终仿真结果尚未冻结，本文在本阶段保留结果表占位，不写入尚未确认的数值结论。

\subsection{5.1 车站背景与模型建立}

龙阳路站连接 2 号线、7 号线、16 号线、18 号线和磁浮线，具有典型的多线换乘属性。站内不同线路的站台、站厅、换乘通道、楼扶梯、闸机和出口共同构成复杂的疏散网络。换乘客流在该站中具有重要机制意义：其路径跨越不同线路空间，需要使用多个共享设施，并会与站台出站流、列车到达流和站厅滞留流在局部设施处叠加。

案例模型根据 CAD 图纸和仿真配置建立站内拓扑。空间建模包括站台、站厅、换乘节点、楼扶梯、闸机、出口前区域和地面出口；设施建模包括有效宽度、服务率、通行方向和共享资源映射；客流建模包括不同线路站台等待人数、列车到达释放人数、站厅滞留人数和换乘关系人数。设施几何和容量参数在最终稿中需要区分 CAD 或模型配置值、文献取值和必要假设值；没有核验的参数不写成官方数据。

\begin{table}[htbp]
\centering
\caption{案例模型的主要组成}
\small
\begin{tabular}{p{0.22\textwidth}p{0.34\textwidth}p{0.32\textwidth}}
\toprule
组成部分 & 建模内容 & 作用 \\
\midrule
空间网络 & 站台、站厅、换乘通道、楼扶梯、闸机和出口节点及其连接边 & 表达乘客从不同来源到出口的可行移动关系 \\
共享设施 & 共用楼扶梯、换乘通道、闸机组和出口前区域的资源映射与服务率 & 表达多来源客流对同一设施容量的竞争 \\
客流来源 & 站台等待、列车到达、站厅滞留和换乘关系批次 & 保留多线路释放和换乘客流的时间结构 \\
评价输出 & 疏散时间、队列等待、来源组清空、出口利用和 Pathfinder 对照指标 & 支撑整体效果和机制解释 \\
\bottomrule
\end{tabular}
\end{table}

\subsection{5.2 疏散场景构造}

实验设置低负荷常规应急场景和高负荷满载列车叠加场景。低负荷常规应急场景包含各线路站台等待、站厅滞留和基础换乘客流，输入总人数为 2187 人。该场景用于检验当共享设施竞争较弱时，AA 是否保持稳定，是否会因过度绕行造成不必要的距离增加或出口偏置。高负荷满载列车叠加场景在基础客流上加入多线路列车到达释放，输入总人数为 17905 人。该场景用于放大多线路站台客流、列车到达客流和换乘客流在共享设施上的同步汇聚效应。

\begin{table}[htbp]
\centering
\caption{低负荷与高负荷疏散场景设置}
\small
\begin{tabular}{p{0.20\textwidth}p{0.14\textwidth}p{0.27\textwidth}p{0.27\textwidth}}
\toprule
场景 & 输入人数 & 主要客流构成 & 实验目的 \\
\midrule
低负荷常规应急 & 2187 & 站台等待、站厅滞留和基础换乘客流 & 检验设施竞争较弱时 AA 的稳定性和低负荷适用性 \\
高负荷满载列车叠加 & 17905 & 基础客流叠加多线路列车到达客流 & 检验多来源客流同步汇聚时 AA 对队列等待和尾部清空的影响 \\
\bottomrule
\end{tabular}
\end{table}

\subsection{5.3 到达负荷与队列预测结果验证}

案例研究首先分析 AA 的中间预测结果，而不是直接进入最终疏散时间比较。该安排与节点通行时间预测、代理模型预测或控制模型识别类论文的实验组织一致：只有先说明模型内部预测是否有效，后续的优化效果才具有可解释性。本文选取换乘通道、楼扶梯、闸机和出口前区域中的关键共享设施，比较预测到达人数、预测队列、实际队列、实际通过人数和未释放人数。

对代表性来源组批次，进一步报告候选路径代价分解。若 AA 改变了某一批次路径，应说明变化来自物理距离、设施等待、批次自身服务时间还是空间回溢等待。若预测队列与执行层实际队列在峰值时段和瓶颈位置上具有相近趋势，则说明到达时刻负荷预测能够捕捉共享设施竞争；若差异较大，则需要从批次粒度、服务率取整、下游接收限制或 Pathfinder 行为规则差异解释。

\begin{table}[htbp]
\centering
\caption{到达负荷与队列预测验证指标}
\small
\begin{tabular}{p{0.22\textwidth}p{0.34\textwidth}p{0.32\textwidth}}
\toprule
分析对象 & 指标 & 说明 \\
\midrule
关键楼扶梯与换乘通道 & 预测到达人数、预测队列、实际队列、峰值等待、最后排队时间 & 判断多来源客流是否在垂向设施和换乘连接处集中汇聚 \\
闸机与出口前区域 & 通过人数、队列人秒、出口分配、来源组构成 & 判断出口方向选择是否造成局部服务压力 \\
代表性来源组批次 & 路径代价分解、候选路径排序、路径重评估记录 & 解释路径改变来自行走距离、设施等待还是空间回溢 \\
\bottomrule
\end{tabular}
\end{table}

\subsection{5.4 整体疏散效率与瓶颈改善分析}

在中间预测结果得到说明后，进一步比较 AA 与 ImprovedAStar 的整体疏散效果。主要指标包括 \(T_{50}\)、\(T_{80}\)、\(T_{95}\)、\(T_{100}\)、平均疏散时间、累计停滞人秒、资源排队人秒、空间阻塞人秒、总移动距离和计算耗时。低负荷场景重点观察 AA 是否引入额外代价；高负荷场景重点观察尾部清空时间、停滞等待和关键设施排队是否变化。

\begin{table}[htbp]
\centering
\caption{低负荷与高负荷场景下的整体疏散性能记录表}
\small
\begin{tabular}{p{0.18\textwidth}p{0.16\textwidth}p{0.10\textwidth}p{0.12\textwidth}p{0.12\textwidth}p{0.14\textwidth}p{0.12\textwidth}}
\toprule
场景 & 方法 & 人数 & \(T_{95}\) & \(T_{100}\) & 累计停滞人秒 & 总移动距离 \\
\midrule
低负荷常规应急 & ImprovedAStar & 2187 & 待冻结 & 待冻结 & 待冻结 & 待冻结 \\
低负荷常规应急 & AA & 2187 & 待冻结 & 待冻结 & 待冻结 & 待冻结 \\
高负荷满载列车叠加 & ImprovedAStar & 17905 & 待冻结 & 待冻结 & 待冻结 & 待冻结 \\
高负荷满载列车叠加 & AA & 17905 & 待冻结 & 待冻结 & 待冻结 & 待冻结 \\
\bottomrule
\end{tabular}
\end{table}

整体指标只回答“是否改善”，不能单独回答“为什么改善”。因此，本节还需要报告关键设施层面的瓶颈变化，包括峰值队列、队列人秒、最后排队时间和设施利用率。若 AA 缩短尾部疏散时间，需要进一步说明这种变化是否来自某些楼扶梯、换乘通道、闸机或出口前区域的持续排队减少。若 AA 降低等待但增加移动距离或计算耗时，也应作为路径组织的代价一并报告。

\subsection{5.5 换乘客流汇聚对整体疏散的作用机制}

换乘客流不是本文唯一评价对象，但它是解释龙阳路站整体疏散效率的重要机制变量。为避免把结果写成单纯算法对比，本文在整体性能分析之后单独分解换乘客流在关键设施上的贡献。对每个关键共享设施，统计站台等待、列车到达、站厅滞留和换乘来源的通过量、等待人秒和峰值队列贡献。若某一设施的排队主要由多个来源组叠加形成，则说明该设施是全站疏散组织中的共享限制环节。

对换乘客流进一步统计其来源线路、主要路径链、关键设施通过量、出口选择和清空时间。该分析回答三个问题：第一，换乘客流是否参与形成关键设施瓶颈；第二，AA 是否改变了换乘客流到达关键设施的时间分布；第三，这种到达分布变化是否进一步影响全站整体疏散效率。通过这一层分析，换乘客流被用于解释整体疏散效率，而不是被写成独立于全站疏散之外的单一人群结果。

\subsection{5.6 Pathfinder 微观仿真对照}

Pathfinder 对照用于检验宏观网络模型识别的拥堵位置和路径组织趋势是否能在微观行人仿真中复现。本文将宏观模型导出的来源组、路径链和路径人数转换为 Pathfinder 中的行人组和行为路径，在相同车站几何、出口设置和客流规模下运行微观仿真。对照指标包括累计疏散曲线、尾部清空时间、拥堵时间、关键设施附近空间拥挤分布和出口利用。

该对照不把 Pathfinder 视为真实车站观测，也不以 Pathfinder 结果替代宏观网络模型结果。其作用是检查 AA 识别出的瓶颈位置是否在微观轨迹中出现相近拥堵，路径组织引起的出口和设施使用变化是否在微观模型中保持同向趋势。若两者存在差异，应从导航网格、个体避让、门选择规则、行为路径约束和容量映射差异解释，而不是简单判定某一模型错误。

\subsection{5.7 适用边界与敏感性实验}

适用边界分析围绕换乘客流比例、列车到达重叠程度和关键共享设施服务能力三类变量展开。选择这三类变量的原因在于，它们直接对应本文的问题机制：换乘客流比例决定换乘流在共享设施竞争中的参与程度，列车到达重叠程度决定多线路客流是否同步释放，关键共享设施服务能力决定瓶颈强度。相比之下，速度-密度阈值主要影响行走时间计算，不是本文队列感知路径组织机制的核心边界。

\begin{table}[htbp]
\centering
\caption{适用边界分析的变量设置}
\small
\begin{tabular}{p{0.22\textwidth}p{0.28\textwidth}p{0.38\textwidth}}
\toprule
变量类别 & 设置方式 & 分析目的 \\
\midrule
换乘客流比例 & 在总客流规模基本一致的前提下调整换乘来源占比 & 判断换乘客流汇聚是否是 AA 产生差异的重要场景条件 \\
列车到达重叠程度 & 调整多线路列车客流释放的时间间隔 & 判断同步到达和错峰到达对共享设施排队和尾部清空的影响 \\
关键共享设施服务能力 & 对主要楼扶梯、换乘通道或闸机服务率进行扰动 & 判断路径组织效果对瓶颈强度的依赖程度 \\
\bottomrule
\end{tabular}
\end{table}

最终结果冻结后，本章结果应按以下顺序呈现。首先报告到达负荷与队列预测结果，证明 AA 的中间机制确实参与路径评价；其次比较低负荷和高负荷场景下 AA 与 ImprovedAStar 的整体疏散性能；再次通过关键设施和换乘客流来源分解解释整体差异；最后通过 Pathfinder 对照和适用边界分析说明结论的跨模型一致性和适用范围。

\end{document}
